\section*{Formalization Tips And Tricks}

\begin{remark*}
When a proof splits an if-and-only-if statement, it is easy to work in the
wrong branch by accident. If the goal is
\[
  A \leftrightarrow B,
\]
then the Lean command
\[
  \texttt{constructor}
\]
creates two proof goals in order:
\[
  A \to B
\]
and then
\[
  B \to A.
\]
Thus a comment such as \texttt{-- =>} labels the forward direction, from the
left side of the iff to the right side, while \texttt{-- <=} labels the reverse
direction, from the right side back to the left side.

For example, in the Euclidean metric proof, the goal has the form
\[
  \texttt{Real.sqrt (...) = 0 <-> p = q}.
\]
After \texttt{constructor}, the first branch is
\[
  \texttt{Real.sqrt (...) = 0 -> p = q},
\]
and the second branch is
\[
  \texttt{p = q -> Real.sqrt (...) = 0}.
\]
So the proof that starts by assuming \texttt{p = q} belongs in the second
branch, not the first.

A useful debugging command is
\[
  \texttt{trace\_state}.
\]
Temporarily placing it inside a proof asks Lean to print the current local
context and goal. For example:
\[
\begin{gathered}
  \texttt{distance\_eq\_zero\_iff\_equal := by}\\
  \texttt{\ \ intro p q}\\
  \texttt{\ \ constructor}\\
  \texttt{\ \ -- =>}\\
  \texttt{\ \ \cdot\ intro hypothesis}\\
  \texttt{\ \ \ \ trace\_state}\\
  \texttt{\ \ \ \ sorry}\\
  \texttt{\ \ -- <=}\\
  \texttt{\ \ \cdot\ intro hypothesis}\\
  \texttt{\ \ \ \ trace\_state}\\
  \texttt{\ \ \ \ sorry.}
\end{gathered}
\]
In the first branch, \texttt{trace\_state} will show that
\texttt{hypothesis} is the assumption that the distance is zero. In the second
branch, it will show that \texttt{hypothesis} is the assumption
\texttt{p = q}. This is often the fastest way to recover your footing when a
proof attempt that feels mathematically right is failing in Lean.

The command is for exploration, not finished code. Once the proof state has
taught you what the branch contains, remove \texttt{trace\_state} and continue
with the actual proof.
\end{remark*}
